% !TeX root = ../../Book5.tex
% 纲要标题：### 9.4 箭图表示的同调初步
% 纲要位置：工作记录/计划纲要.md，第 4029 行。

\section{箭图表示的同调初步}
\label{b5:sec:0904}


% 纲要范围：
%
% * 无环箭图的投射表示
% * 短投射消解与 Ext^1
% * 第10章根系分类的准备
% * 本节的 \(\operatorname{Ext}^1\) 可先按短正合扩张定义，并就地说明与短投射消解计算的一致性；不要求完整同调代数课程
%

区间分解说明线性定向箭图很容易分类，但一般箭图不能同时把所有箭化为这种形式。另一种有效做法是把一个表示写成若干简单的投射表示之间的商。对有限无环箭图，其无关系路径代数 \(kQ\) 的每个有限维表示都有长度至多一的投射消解，因而可以用一条短消解同时计算同态与扩张。

\subsection{无环箭图的投射表示}
\label{b5:subsec:0904:01}

设 \(k\) 为域，\(Q\) 为有限箭图，\(P\) 为 \(Q\) 的有限维 \(k\) 表示。若对有限维 \(k\) 表示之间的每个满射同态 \(E\to V\)，任意表示同态 \(P\to V\) 都能提升到 \(P\to E\)，则称 \(P\) 为投射表示。这里满射按顶点理解，定义的是有限维表示范畴中的投射性。

\ProseParagraphBreak

下面设 \(Q\) 无环。由\cref{b5:prop:finite-path-algebra}，路径总数有限，\(A=kQ\) 以及各个 \(Ae_i\) 都有限维。此时上述投射性与通常的左 \(A\)-模投射性一致。事实上，任一有限维左 \(A\)-模 \(P\) 都有满射 \(A^n\to P\)，其中 \(n\) 有限；\(A^n\) 也是有限维对象。若 \(P\) 在有限维模范畴中投射，该满射便有截面，使 \(P\) 成为有限自由模的直和因子，因而对所有左 \(A\)-模之间的满射都满足提升性质。反向则直接限制到有限维对象。
\begin{proposition}\label{b5:prop:vertex-projective}
设 \(k\) 为域，\(Q\) 为有限无环箭图，\(A=kQ\)，\(i\in Q_0\)，\(e_i\) 为顶点 \(i\) 的长度零路径。左 \(A\)-模 \(P_i=Ae_i\) 在顶点 \(j\) 以全部路径 \(i\to j\) 为基，箭通过在路径末端继续拼接而作用。对 \(Q\) 的任意有限维 \(k\) 表示 \(W\)，在 \(e_i\) 处求值给出自然同构
\[
\operatorname{Hom}_Q(P_i,W)\cong W_i,\qquad f\longmapsto f(e_i).
\]
因此 \(P_i\) 投射，且 \(\operatorname{End}_Q(P_i)\cong k\)，所以不可分解。
\end{proposition}
\symidx{\(P_i=kQe_i\)：无环箭图的顶点投射表示}
\begin{proof}
模同态满足 \(f(e_i)=e_if(e_i)\in W_i\)，又满足 \(f(p)=W_p f(e_i)\)。反过来，任给 \(w\in W_i\)，这个公式在路径基上线性延拓，便给出唯一模同态。该求值映射对 \(W\) 的同态相容，因而自然。顶点取值函子保持正合列，特别是满射，所以该 Hom 函子也保持满射，得到投射性。无环性使 \(i\to i\) 只有长度零路径；取 \(W=P_i\)，自同态由 \(e_i\) 的标量倍决定。其自同态环是域，因而无非平凡直和投影。
\end{proof}
在 \(Q=(1\xrightarrow{a}2\xrightarrow{b}3)\) 中，\(P_1=kQe_1\) 在三个顶点上分别以 \(e_1,a,ba\) 为基；两支箭把 \(e_1\) 送到 \(a\)，再把 \(a\) 送到 \(ba\)。取目标表示
\[
W_1\xrightarrow{f}W_2\xrightarrow{g}W_3.
\]
指定 \(\Phi_1(e_1)=w\in W_1\) 后，箭的交换条件要求
\[
\Phi_2(a)=f(w),\qquad \Phi_3(ba)=gf(w).
\]
这些公式对任意 \(w\) 都给出唯一的表示同态 \(\Phi:P_1\to W\)，所以求值同构确实取遍整个 \(W_1\)。若换成简单顶点模 \(S_1=(k\to0\to0)\)，第二个顶点的同态分量必为零，第一支箭的交换条件便要求 \(f(w)=0\)，故
\[
\operatorname{Hom}_Q(S_1,W)\cong\ker f.
\]
例如 \(W_1\ne0\) 且 \(f\) 单射时，后一个 Hom 空间为零，而 \(\operatorname{Hom}_Q(P_1,W)\) 非零。\(P_1\) 允许一个生成元沿所有路径自由传播，\(S_1\) 则把正长度路径上的传播规定为零。沿满射提升 \(w\) 后，各条路径上的像就由上述公式自动确定，这也说明了顶点投射模的提升性质。

\ProseParagraphBreak

对 \(1\to2\to\cdots\to n\)，有 \(P_i=I[i,n]\)。箭向改变时投射表示也改变：在 \(1\to2\leftarrow3\) 中，\(P_2\) 只支集在顶点2，\(P_1,P_3\) 则各沿一支箭延伸。判断投射性应看“从这个顶点出发的路径”，不是只凭支集大小。

\subsection{\texorpdfstring{短投射消解与 \(\operatorname{Ext}^1\)}{短投射消解与 Ext1}}
\label{b5:subsec:0904:02}

先看单箭 \(1\xrightarrow{a}2\) 上的表示 \(V_1\xrightarrow{f}V_2\)。中项 \(P_1\otimes_kV_1\oplus P_2\otimes_kV_2\) 在顶点 \(2\) 是 \(V_1\oplus V_2\)：第一份来自沿箭自由传播的 \(V_1\)，第二份是原来放在顶点 \(2\) 的向量。因此下面的消解在顶点 \(2\) 具体成为
\[
0\longrightarrow V_1\xrightarrow{v\mapsto(v,-f(v))}
V_1\oplus V_2\xrightarrow{(x,y)\mapsto f(x)+y}
V_2\longrightarrow0.
\]
在顶点 \(1\)，它是 \(0\to0\to V_1\xrightarrow{1}V_1\to0\)。向量 \(v\) 自由传播后的像与原表示中的 \(f(v)\) 必须相等，两者的差 \((v,-f(v))\) 正是要模掉的关系。一般箭图中，每支箭都产生同样的一组关系。

\ProseParagraphBreak

再取两支连续箭 \(1\xrightarrow{a}2\xrightarrow{b}3\)，以及表示 \(V_1\xrightarrow{f}V_2\xrightarrow{g}V_3\)。中项在顶点 \(3\) 是 \(V_1\oplus V_2\oplus V_3\)：三份向量分别从顶点 \(1\) 沿 \(ba\)、从顶点 \(2\) 沿 \(b\)，以及从顶点 \(3\) 沿 \(e_3\) 到达这里。因此增广映射为
\[
\epsilon_3(x,y,z)=gf(x)+g(y)+z.
\]
左端在顶点 \(3\) 是 \(V_1\oplus V_2\)。箭 \(a\) 的关系沿 \(b\) 传播后为 \((u,-f(u),0)\)，箭 \(b\) 的关系为 \((0,v,-g(v))\)，合在一起给出
\[
\mathrm{d}_3(u,v)=(u,-f(u)+v,-g(v)).
\]
若 \(\mathrm{d}_3(u,v)=0\)，第一坐标使 \(u=0\)，第二坐标再使 \(v=0\)，所以它单射。直接代入有
\[
\epsilon_3\mathrm{d}_3(u,v)
=gf(u)+g(-f(u)+v)-g(v)=0.
\]
反过来，若 \(\epsilon_3(x,y,z)=0\)，取 \(u=x\)、\(v=y+f(x)\)，则 \(-g(v)=z\)，所以 \((x,y,z)=\mathrm{d}_3(u,v)\)。\(\epsilon_3\) 又因 \(\epsilon_3(0,0,z)=z\) 为满射，故这一顶点处的序列正合。

\ProseParagraphBreak

长度二路径并不需要再添独立的传播关系。在模掉上述两支箭的关系以后，有
\[
ba\otimes u=b\otimes f(u)=e_3\otimes gf(u).
\]
第一步用箭 \(a\) 的关系并沿 \(b\) 继续传播，第二步用箭 \(b\) 的关系。更长路径也逐支箭作同样的替换，因此全部路径的传播关系都由逐箭关系及其继续拼接生成。这正是消解左端按箭索引的原因。

\begin{theorem}\label{b5:thm:acyclic-quiver-resolution}
设 \(k\) 为域，\(Q\) 为有限无环箭图，\(A=kQ\)，\(V=(V_i,V_a)\) 为 \(Q\) 的有限维 \(k\) 表示。记 \(e_i\) 为顶点 \(i\) 的长度零路径、\(P_i=Ae_i\)，\(V_p\) 为沿路径 \(p\) 的箭映射复合。存在左 \(A\)-模的短正合列
\[
0\longrightarrow
\bigoplus_{a:i\to j}P_j\otimes_k V_i
\xrightarrow{\mathrm{d}}
\bigoplus_{i\in Q_0}P_i\otimes_k V_i
\xrightarrow{\epsilon}V\longrightarrow0.
\]
其中左 \(A\) 作用在每项的第一个因子上，\(\epsilon(p\otimes v)=V_pv\)。对于属于箭 \(a:i\to j\) 的左端分量，微分为
\[
\mathrm{d}(p\otimes v)=pa\otimes v-p\otimes V_av,
\]
第一项放在中项的 \(i\) 分量，第二项放在 \(j\) 分量。两项直和都是有限维投射表示。
\end{theorem}
\begin{proof}
\(\epsilon\) 与 \(\mathrm{d}\) 都为 \(A\) 线性，且 \(\epsilon\mathrm{d}=0\)，因为路径 \(pa\) 的作用是先 \(V_a\) 再 \(V_p\)。每个顶点的向量由 \(e_i\otimes v\) 提升，所以 \(\epsilon\) 满。
证明 \(\mathrm{d}\) 单射时，对左端的一个非零元素按路径 \(p\) 收集其向量系数，选出现的最大路径长度 \(L\)。正项 \(pa\otimes v\) 长度为 \(L+1\)，负项长度至多 \(L\)。一条正长度路径唯一决定其第一支箭 \(a\) 及剩余路径 \(p\)，所以这些最高长度正项之间也不能相消。因此像非零。
在中项模掉 \(\operatorname{im}\mathrm{d}\) 后，关系
\(pa\otimes v=p\otimes V_av\)
可逐次删去路径的第一支箭，把 \(p\otimes v\) 化为 \(e_{t(p)}\otimes V_pv\)。于是每个商类都是长度零项的和；若其 \(\epsilon\) 像为零，各顶点的向量系数都为零，商类也为零。故 \(\ker\epsilon=\operatorname{im}\mathrm{d}\)。最后由\cref{b5:prop:vertex-projective}，各项是有限个 \(P_i\) 的直和，因而投射。
\end{proof}
无环性在这里保证各个 \(P_i\) 有限维；顶点、箭以及 \(V_i\) 的维数也都有限，所以消解的两项仍在有限维表示范畴中。有向圈本身并不意味着长度一投射消解不存在。例如单圈的路径代数为 \(k[t]\)，对任意 \(\lambda\in k\)，一维模 \(k[t]/(t-\lambda)\) 有消解
\[
0\longrightarrow k[t]\xrightarrow{\times(t-\lambda)}k[t]
\longrightarrow k[t]/(t-\lambda)\longrightarrow0.
\]
首个映射单射，因为 \(k[t]\) 是整环且 \(t-\lambda\ne0\)；其像为 \((t-\lambda)\)，恰是商映射的核。两个 \(k[t]\) 都是秩一自由模，故这是一条长度一投射消解，但它们在 \(k\) 上无限维，属于全体左模的范畴。加入路径关系则是另一种变化：关系会把某些原来的非零路径规定为零，可能使上述逐箭映射产生新的核。

\ProseParagraphBreak

为了说明标准投射消解怎样计算扩张，先把本节所用的 \(\operatorname{Ext}^1\) 定义清楚。对有限维 \(A\) 模 \(V,W\)，考虑短正合列
\[
0\longrightarrow W\longrightarrow E\longrightarrow V\longrightarrow0.
\]
两列若有一个中项同构使两端恒等及相应图表交换，就看成同一个扩张类，记全体类为 \(\operatorname{Ext}_A^1(V,W)\)。零类是分裂列。同样的定义也适用于一般左模；本章的主要应用仍取有限维对象。对两个扩张先取直和，得到核为 \(W\oplus W\)、商为 \(V\oplus V\) 的短正合列，再沿对角映射 \(V\to V\oplus V\)、\(v\mapsto(v,v)\) 拉回。其新中项由两份中项中商像相等的元素对组成，仍满射到 \(V\)，核为 \(W\oplus W\)。再沿相加映射 \(W\oplus W\to W\)、\((w_1,w_2)\mapsto w_1+w_2\) 推出，就得到核为 \(W\)、商为 \(V\) 的新扩张，定义两个类的和。沿标量映射 \(W\to W\) 推出定义标量乘法。拉回与推出保持这里的短正合列，且两份扩张中项的等价同构分别诱导新中项的同构，所以这些运算不依赖代表列。下面的具体描述再证明它们满足 \(k\) 向量空间的运算规律。
\begin{proposition}\label{b5:prop:projective-presentation-extensions}
设 \(k\) 为域，\(A\) 为含幺 \(k\) 代数，\(V,W\) 为左 \(A\)-模，并给定左 \(A\)-模的短正合列
\(0\to K\xrightarrow{\iota}P\to V\to0\)，其中 \(P\) 投射。对于固定的这条短正合列，有关于 \(W\) 自然的 \(k\) 线性同构
\[
\operatorname{Ext}_A^1(V,W)\cong
\operatorname{coker}\left(
\operatorname{Hom}_A(P,W)\xrightarrow{\iota^*}
\operatorname{Hom}_A(K,W)\right).
\]
这里 \(\iota^*\) 把同态 \(f:P\to W\) 送到 \(f\circ\iota:K\to W\)，\(\operatorname{Ext}_A^1(V,W)\) 取短正合扩张类及其 Baer 和。
\end{proposition}
\begin{proof}
给定 \(f:K\to W\)，映射 \((f,-\iota):K\to W\oplus P\)、\(x\mapsto(f(x),-\iota(x))\) 为 \(A\) 线性，其像是 \(A\)-子模。取
\[
E_f=(W\oplus P)/\operatorname{im}(f,-\iota).
\]
映射 \(w\mapsto\bigl[(w,0)\bigr]\) 单射，商到 \(V\) 的映射为 \(\bigl[(w,p)\bigr]\mapsto\bar p\)，其核正是这份 \(W\)，所以得到扩张。反过来，对任一扩张，利用 \(P\) 的投射性将 \(P\to V\) 提升为 \(u:P\to E\)。记原扩张的商映射为 \(\pi:E\to V\)、核嵌入为 \(j:W\to E\)。提升条件使 \(\pi u\iota=0\)，故 \(u\iota\) 落在 \(\ker\pi=j(W)\)，唯一给出 \(f:K\to W\) 使 \(u\iota=jf\)。映射 \((w,p)\mapsto j(w)+u(p)\) 零化每个 \((f(x),-\iota(x))\)，因而在 \(E_f\) 上良定义，并在核 \(W\) 与商 \(V\) 上均为恒等。它满射，因为任意 \(e\in E\) 的商像都可由某个 \(p\in P\) 提升，余差 \(e-u(p)\) 属于 \(j(W)\)。若 \(j(w)+u(p)=0\)，则 \(p=\iota(x)\) 对某个 \(x\in K\) 成立，继而 \(w=-f(x)\)，所以 \((w,p)\) 在商去的子模中；这证明单射，得到 \(E_f\cong E\)。
若 \(u,u'\) 是两个提升，则 \(\pi(u'-u)=0\)，所以它们的差唯一写成 \(jh\)，其中 \(h:P\to W\)。于是对应的限制映射满足 \(f'-f=h\iota\)，只改变一个来自 \(\operatorname{Hom}_A(P,W)\) 的项。反过来，若两个 \(f\) 的差可延伸到 \(P\)，用这个延伸修正 \(W\oplus P\) 的第一坐标，便得到固定两端的扩张同构。因此商空间恰好记录扩张类。按定义作拉回和推出时，两个限制映射相加，标量推出则把限制映射乘相应标量；所以对应保持上述运算。

对任意 \(A\)-模同态 \(h:W\to W'\)，映射
\[
E_f\longrightarrow E_{h\circ f},\qquad
[(w,p)]\longmapsto[(h(w),p)]
\]
把子模上的映射变为 \(h\)，在商模 \(V\) 上诱导恒等；由商模的构造，\(E_{h\circ f}\) 正是沿 \(h\) 推出所得扩张。另一方面，余核上的诱导映射也把 \([f]\) 送到 \([h\circ f]\)。故上述同构关于 \(W\) 自然。
\end{proof}

\subsection{同调方法与箭图表示的分类}
\label{b5:subsec:0904:03}

\begin{theorem}\label{b5:thm:quiver-hom-ext-euler}
设 \(k\) 为域，\(Q\) 为有限无环箭图，\(V=(V_i,V_a)\)、\(W=(W_i,W_a)\) 为有限维 \(k\) 表示。记 \(\operatorname{Ext}_Q^1(V,W)=\operatorname{Ext}_{kQ}^1(V,W)\)，其中两表示视作左 \(kQ\)-模。存在自然正合列
\[
\begin{split}
0\longrightarrow\operatorname{Hom}_Q(V,W)
&\longrightarrow\bigoplus_{i\in Q_0}\operatorname{Hom}_k(V_i,W_i)\\
&\xrightarrow{\delta}
\bigoplus_{a:i\to j}\operatorname{Hom}_k(V_i,W_j)
\longrightarrow\operatorname{Ext}_Q^1(V,W)\longrightarrow0,
\end{split}
\]
其中 \(\delta\bigl((f_i)_i\bigr)=(W_af_i-f_jV_a)_{a:i\to j}\)。若 \(d_i=\dim_k V_i\)、\(e_i=\dim_k W_i\)，则
\[
\dim_k\operatorname{Hom}_Q(V,W)-\dim_k\operatorname{Ext}_Q^1(V,W)
=\sum_i d_i e_i-\sum_{a:i\to j}d_i e_j.
\]
\end{theorem}
\symidx{\(\operatorname{Hom}_Q(V,W),\operatorname{Ext}_Q^1(V,W)\)：箭图表示的同态空间与短正合扩张类群}
\begin{proof}
将\cref{b5:prop:projective-presentation-extensions}用于\cref{b5:thm:acyclic-quiver-resolution}。由\cref{b5:prop:vertex-projective}及张量的线性泛性质，
\[
\operatorname{Hom}_Q(P_i\otimes_k V_i,W)
\cong\operatorname{Hom}_k(V_i,W_i).
\]
将同态应用于 \(\mathrm{d}(p\otimes v)\)，再令 \(p=e_j\)，便得到 \(W_af_i-f_jV_a\)。这个映射的核正是表示同态的交换条件，而余核由前一命题计算扩张，所以所列序列正合。所有空间有限维，维数的交错和为零，得到公式。
\end{proof}
若 \(S_i\) 表示只在顶点 \(i\) 放置 \(k\) 的表示，则对于不同顶点 \(i,j\)，有
\[
\dim\operatorname{Ext}_Q^1(S_i,S_j)
=\#\{\,a:i\to j\,\}.
\]
这里右侧的方向与左模约定一致：在 \(1\to2\) 中，非分裂列是
\(0\to S_2\to I[1,2]\to S_1\to0\)。
\ProseParagraphBreak
% 跨卷引用目标：b4:thm:ext-balanced（第四卷定理9.2.2）。
令 \(V=W\)，公式左侧比较自同态与自扩张，右侧只含维数向量。下一章将把这个整数二次式与反射联系起来。由第四卷定理9.2.2的 Ext 平衡性，高次 Ext 可以用投射消解经 \(\operatorname{Hom}_{kQ}(-,W)\) 后的上同调。上面的消解只有次数零与一的投射项，所得 Hom 复形在次数 \(r\ge2\) 处为零，因此 \(\operatorname{Ext}_{kQ}^r(V,W)=0\) 对所有 \(r\ge2\) 成立。加上关系以后，长度一消解一般不再成立。

\subsection{短正合扩张与消解计算的一致性}
\label{b5:subsec:0904:04}

上述余核还能从扩张中的箭矩阵直接读出。给定
\(0\to W\to E\to V\to0\)，在每个顶点选择向量空间截面，识别
\(E_i=W_i\oplus V_i\)。箭矩阵必为
\[
E_a=
\begin{pmatrix}
W_a&\xi_a\\0&V_a
\end{pmatrix},
\qquad \xi_a:V_i\to W_j\quad(a:i\to j).
\]
没有关系时，各个 \(\xi_a\) 可以任意选择。将截面 \(s_i\) 换成
\(s_i+f_i\)、\(f_i:V_i\to W_i\)，右上角变为
\[
\xi'_a=\xi_a+W_af_i-f_jV_a.
\]
这是将 \(E_a\bigl(s_i(v)+f_i(v)\bigr)\) 减去 \(s_j(V_av)+f_j(V_av)\) 的直接计算。因此两组矩阵表示相同的扩张类，当且仅当相差 \(\delta\bigl((f_i)_i\bigr)\)；中项同构固定两端时也恰有这种分块上三角形式。
\ProseParagraphBreak
这与消解所得的对应一致，而不只是维数相同。由各顶点的截面构造中项到 \(E\) 的提升
\[
P_i\otimes V_i\longrightarrow E,\qquad p\otimes v\longmapsto E_p s_i(v).
\]
在箭 \(a:i\to j\) 所对应的关系 \(a\otimes v-e_j\otimes V_av\) 上，它的值为
\(E_as_i(v)-s_j(V_av)=\xi_a(v)\)，正是前一命题取到的限制映射。
\begin{example}\label{b5:exm:arrow-extension-parameter}
设 \(k\) 为域，\(Q=(1\to2)\)。扩张 \(0\to S_2\to E\to S_1\to0\) 由一个标量 \(c\) 给出，中项为 \(E_c=(k\xrightarrow{c}k)\)，子模 \(S_2\) 在顶点 \(2\) 的嵌入及商模 \(S_1\) 在顶点 \(1\) 的商映射都取恒等。此时两个顶点的 \(\operatorname{Hom}_k(V_i,W_i)\) 都为零，所以不同 \(c\) 给出不同扩张类，\(\operatorname{Ext}_Q^1(S_1,S_2)\cong k\)。

表示同构 \(E_c\to E_{c'}\) 的两个顶点分量是非零标量 \(u,v\)，箭的交换条件为
\[
vc=c'u.
\]
若要求它给出固定两端的扩张等价，顶点 \(1\) 的商模上恒等强制 \(u=1\)，顶点 \(2\) 的子模上恒等强制 \(v=1\)，于是 \(c'=c\)。若只比较中项，则 \(u,v\) 可以任取非零值，参数相应变为 \(c'=\dfrac vu c\)。这一换基作用只有两个轨道 \(\{0\}\) 与 \(k^\times\)：\(c=0\) 时中项为 \(S_1\oplus S_2\)，而 \(c,c'\ne0\) 时取 \(u=1\)、\(v=\dfrac{c'}c\) 即得同构，所以所有非零参数的中项都同构于 \(I[1,2]\)。扩张类记录固定两端的短正合列，而中项同构还允许在两端换基，因此两种分类给出的参数不同。
\end{example}
对于有关系箭图，还须把分块矩阵代入关系。比如 \(1\xrightarrow{a}2\xrightarrow{b}3\) 加上 \(ba=0\) 时，右上角应满足
\(W_b\xi_a+\xi_bV_a=0\)，并非每组 \(\xi_a,\xi_b\) 都可用。这里已经能看见关系怎样限制扩张；不能继续把无关系箭图的 Euler 公式\footnote{欧拉（Leonhard Euler，1707-1783），瑞士数学家，研究数论、分析与图论，奠定多项基础。}原样照搬。
